\chapter{PENDAHULUAN}
Pada bab ini dijelaskan tentang latar belakang, rumusan masalah, batasan masalah, tujuan penelitian, dan manfaat penelitian.

\section{Latar Belakang}
Uraikan konteks penelitian, untuk apa penelitian ini dilakukan, dan hal apa yang mengarahkan penelitian ini. Latar belakang memuat teori utama yang relevan dan terbaru \parencite{contoh2024artikel} yang dipadukan sehingga mengerucut pada suatu persoalan yang kemudian disusun dalam bentuk rumusan masalah. Sitasi naratif ditulis dengan \texttt{\textbackslash textcite}, misalnya \textcite{penulis2015buku}.

\section{Rumusan Masalah}
Rumusan masalah dalam penelitian ini adalah sebagai berikut:
\begin{enumerate}
  \item Rumusan masalah pertama?
  \item Rumusan masalah kedua?
\end{enumerate}

\section{Batasan Masalah}
Batasan masalah dalam penelitian ini adalah sebagai berikut:
\begin{enumerate}
  \item Batasan pertama.
  \item Batasan kedua.
\end{enumerate}

\section{Tujuan Penelitian}
Berdasarkan rumusan masalah, tujuan dari penelitian ini adalah sebagai berikut:
\begin{enumerate}
  \item Tujuan yang menjawab rumusan masalah pertama.
  \item Tujuan yang menjawab rumusan masalah kedua.
\end{enumerate}

\section{Manfaat Penelitian}
Manfaat yang diharapkan dari penelitian ini adalah sebagai berikut:
\begin{enumerate}
  \item Secara teoritis, \dots
  \item Secara praktis, \dots
  \item Sebagai bahan referensi untuk penelitian selanjutnya.
\end{enumerate}
